\documentclass[a4paper,12pt]{article}

% --- Pakiety ---
\usepackage[polish]{babel}
\usepackage[utf8]{inputenc}
\usepackage[T1]{fontenc}
\usepackage[top=2.5cm,bottom=2.5cm,left=3cm,right=3cm]{geometry}
\usepackage{amsmath}
\usepackage{graphicx}
\usepackage{booktabs}
\usepackage[colorlinks=true, allcolors=blue]{hyperref}

% --- Tytuł ---
\title{
    \vspace{1.5cm}
    \textbf{PODSTAWY TEORII INFORMACJI}\\
    \vspace{1.5cm}
    Laboratorium 2\\
    Analiza stanu wiedzy\\
    \vspace{1.5cm}
    \huge \textbf{Analiza informacji i wiarygodności prognoz meteorologicznych}\\
    \vspace{1cm}
}

\author{
    Maciek Chotkowski \\ Martsin Lazouski \\ Nadia Ziecina \\
    Michał Kasznia \\ Kasiński Julian \\ Volodymyr Kamenchuk \\ Usevalad Nikalaichuk 
}

\date{\today}

\begin{document}

\maketitle
\newpage

\tableofcontents
\newpage

%--------------------------------------------------
\section{Wprowadzenie}

Prognozy meteorologiczne stanowią istotne źródło informacji wykorzystywanej w codziennym podejmowaniu decyzji. Pomimo szerokiej dostępności danych, różne modele numeryczne często generują rozbieżne wyniki, co prowadzi do niepewności użytkownika.

Celem niniejszego etapu jest analiza problemu nie tylko jako zadania numerycznego, lecz przede wszystkim jako problemu informacji oraz jej użyteczności.

%--------------------------------------------------
\section{Informacja w prognozach meteorologicznych}

\subsection{Poziom syntaktyczny}
Na poziomie syntaktycznym prognoza jest zbiorem danych liczbowych, takich jak temperatura, opady czy prędkość wiatru.

\subsection{Poziom semantyczny}
Na poziomie semantycznym dane te nabierają znaczenia, np. temperatura może być interpretowana jako „zimno” lub „ciepło”.

\subsection{Poziom pragmatyczny}
Na poziomie pragmatycznym informacja wpływa na decyzje użytkownika, np. wybór ubioru lub planowanie aktywności.

%--------------------------------------------------
\section{Model użytkownika}

Użytkownik końcowy nie analizuje wartości liczbowych ani miar błędu, lecz oczekuje informacji pozwalającej na podjęcie decyzji.

Najważniejsze dla użytkownika są:
\begin{itemize}
    \item wystąpienie opadów,
    \item przekroczenie progów temperatury,
    \item ogólna wiarygodność prognozy.
\end{itemize}

%--------------------------------------------------
\section{Analiza istniejących rozwiązań}

\subsection{Metryki statystyczne}
Najczęściej stosowane są miary błędu takie jak MAE, MSE oraz RMSE.

\subsection{Podejścia probabilistyczne}
Prognozy probabilistyczne pozwalają uwzględnić niepewność danych.

\subsection{Analiza zdarzeń}
Często analizuje się konkretne zdarzenia, takie jak wystąpienie opadów.

\subsection{Modele prognostyczne}
Do najczęściej stosowanych modeli należą GFS, ECMWF oraz UM.

%--------------------------------------------------
\section{Reprezentacja informacji}

Dane meteorologiczne reprezentowane są jako szeregi czasowe.

Istotnym zagadnieniem jest selekcja informacji, czyli odrzucenie nieistotnych różnic (np. różnicy rzędu 0.5°C).

%--------------------------------------------------
\section{Metryki i kryteria oceny}

Oprócz klasycznych miar błędu należy uwzględnić metryki użytkowe:
\begin{itemize}
    \item trafność przewidywania opadów,
    \item błędy przekroczenia progów,
    \item stabilność prognoz.
\end{itemize}

%--------------------------------------------------
\section{Wyzwania i ograniczenia}

\begin{itemize}
    \item chaotyczny charakter pogody,
    \item błędy pomiarowe,
    \item wpływ warunków lokalnych,
    \item różnice między modelami.
\end{itemize}

%--------------------------------------------------
\section{Wnioski}

Prognoza meteorologiczna nie jest jedynie zbiorem danych, lecz informacją, której wartość zależy od kontekstu oraz użytkownika.

Ocena jakości prognozy powinna uwzględniać nie tylko dokładność numeryczną, ale również jej użyteczność w praktyce.

\end{document}